% Part: normal-modal-logic
% Chapter: tableaux
% Section: countermodels

\documentclass[../../../include/open-logic-section]{subfiles}

\begin{document}

\olfileid{nml}{tab}{cou}

\olsection{从真值树构建反模型}

\begin{explain}
  完备性定理的证明不仅仅表明若 $\Entails !A$ 则 $\Proves !A$，
  它还在 $\Entails/ A$ 的情况下，为我们提供了构造~$!A$ 的\emph{反模型}的方法。
  就~$\Log{K}$ 而言，这一方法构成了一个\emph{判定过程}。假设 $\Entails/ !A$，
  \olref[cpl]{prop:complete-tableau} 的证明便给出了构造一个完全分析表的方法。
  该方法事实上总是会\emph{终止}。$\Log{K}$ 的命题规则只添加复杂度较低的带前缀公式，
  也就是说，对于任何带号公式 $\sFmla{S}{!A}[\sigma]$，每条命题规则在分支上只需应用一次。
  新前缀仅由\iftag{prvBox}{$\TRule{\False}{\Box}$}{}\iftag{notprvBox,notprvDiamond}{}{
    和 }\iftag{prvDiamond}{$\TRule{\True}{\Diamond}$}{}
  \iftag{notprvBox,notprvDiamond}{规则}{这些规则}生成，
  且同样只需应用一次（并只产生一个新前缀）。
  \iftag{prvBox}{$\TRule{\True}{\Box}$}{}\iftag{notprvBox,notprvDiamond}{}{
    和 }\iftag{prvDiamond}{$\TRule{\False}{\Diamond}$}{}
  可能\iftag{notprvBox,notprvDiamond}{需要}{需要}应用多次，
  但每个前缀只需应用一次，且生成的新前缀只有有限多个。
  因此，该构造在经过有限多步后，要么得到一个封闭分支，要么得到一个完全分支。

  一旦构造出包含开放完全分支的分析表，
  \olref[cpl]{thm:tableau-completeness} 的证明就为我们提供了一个满足原始带前缀公式集的显式模型。
  因此，不仅当 $\Gamma \Entails !A$ 时存在封闭分析表且 $\Gamma \Proves !A$；
  而且如果我们以正确的方式寻找封闭分析表并最终得到一个“完全的”分析表，
  我们不仅能知道 $\Gamma \Entails/ !A$，还能实际构造出反模型。
\end{explain}

\iftag{prvBox}{
\begin{ex}
  我们知道 $\Proves/ \Box(p \lor q) \lif (\Box p \lor \Box
  q)$。分析表的构造始于：
  \begin{oltableau}
    [\pFmla{\False}{\Box(p \lor q) \lif (\Box p \lor \Box q)}{1},
      just = \TAss, checked
      [\pFmla{\True}{\Box(p \lor q)}{1},
        just = {\TRule{\False}{\lif}[1]},
        [\pFmla{\False}{\Box p \lor \Box q}{1},
          just = {\TRule{\False}{\lif}[1]}, checked
          [\pFmla{\False}{\Box p}{1},
            just = {\TRule{\False}{\lor}[3]}, checked
            [\pFmla{\False}{\Box q}{1},
              just = {\TRule{\False}{\lor}[3]}, checked
              [\pFmla{\False}{p}{1.1},
                just = {\TRule{\False}{\Box}[4]}, checked
                [\pFmla{\False}{q}{1.2},
                  just = {\TRule{\False}{\Box}[5]}, checked
                ]
              ]
            ]
          ]
        ]
      ]
    ]
  \end{oltableau}
  这个分析表尚未完成。下一步，我们考虑唯一没有对勾标记的行：第~$2$ 行上的带前缀公式
  $\sFmla{\True}{\Box(p \lor q)}[1]$。根据封闭分析表的构造要求，需要对分支上使用的每个前缀应用 $\TRule{\True}{\Box}$ 规则，即对 $1.1$ 和 $1.2$ 都要应用：
  \begin{oltableau}
    [\pFmla{\False}{\Box(p \lor q) \lif (\Box p \lor \Box q)}{1},
      just = \TAss, checked
      [\pFmla{\True}{\Box(p \lor q)}{1},
        just = {\TRule{\False}{\lif}[1]},
        [\pFmla{\False}{\Box p \lor \Box q}{1},
          just = {\TRule{\False}{\lif}[1]}, checked
          [\pFmla{\False}{\Box p}{1},
            just = {\TRule{\False}{\lor}[3]}, checked
            [\pFmla{\False}{\Box q}{1},
              just = {\TRule{\False}{\lor}[3]}, checked
              [\pFmla{\False}{p}{1.1},
                just = {\TRule{\False}{\Box}[4]}, checked
                [\pFmla{\False}{q}{1.2},
                  just = {\TRule{\False}{\Box}[5]}, checked
                  [\pFmla{\True}{p \lor q}{1.1},
                    just = {\TRule{\True}{\Box}[2]}
                    [\pFmla{\True}{p \lor q}{1.2},
                      just = {\TRule{\True}{\Box}[2]}
                    ]
                  ]
                ]
              ]
            ]
          ]
        ]
      ]
    ]
  \end{oltableau}
  现在第 2、8、9 行没有对勾标记。但由于没有添加新的前缀，所以我们对第 8 行和第 9 行应用 $\TRule{\True}{\lor}$ 规则，应用于所有由此产生且尚未闭合的分支：
  \begin{oltableau}
    [\pFmla{\False}{\Box(p \lor q) \lif (\Box p \lor \Box q)}{1},
      just = \TAss, checked
      [\pFmla{\True}{\Box(p \lor q)}{1},
        just = {\TRule{\False}{\lif}[1]}, checked
        [\pFmla{\False}{\Box p \lor \Box q}{1},
          just = {\TRule{\False}{\lif}[1]}, checked
          [\pFmla{\False}{\Box p}{1},
            just = {\TRule{\False}{\lor}[3]}, checked
            [\pFmla{\False}{\Box q}{1},
              just = {\TRule{\False}{\lor}[3]}, checked
              [\pFmla{\False}{p}{1.1},
                just = {\TRule{\False}{\Box}[4]}, checked
                [\pFmla{\False}{q}{1.2},
                  just = {\TRule{\False}{\Box}[5]}, checked
                  [\pFmla{\True}{p \lor q}{1.1},
                    just = {\TRule{\True}{\Box}[2]}, checked
                    [\pFmla{\True}{p \lor q}{1.2},
                      just = {\TRule{\True}{\Box}[2]}, checked
                      [\pFmla{\True}{p}{1.1},
                        just = {\TRule{\True}{\lor}[8]}, checked, close
                      ]
                      [\pFmla{\True}{q}{1.1},
                        just = {\TRule{\True}{\lor}[8]}, checked
                        [\pFmla{\True}{p}{1.2},
                          just = {\TRule{\True}{\lor}[9]}, checked]
                        [\pFmla{\True}{q}{1.2},
                          just = {\TRule{\True}{\lor}[9]}, checked, close]
                      ]
                    ]
                  ]
                ]
              ]
            ]
          ]
        ]
      ]
    ]
  \end{oltableau}
  剩下一个开放分支，并且它是完全的。我们据此定义模型：可能世界集为 $W = \{1, 1.1, 1.2\}$
  （即该开放分支上仅出现的前缀），可达关系为 $R = \{\tuple{1, 1.1}, \tuple{1, 1.2}\}$，赋值为 $V(p) =
  \{1.2\}$（因为第~11 行包含 $\sFmla{\True}{p}[1.2]$）和 $V(q) = \{1.1\}$（因为第~10 行包含 $\sFmla{\True}{q}[1.1]$）。该模型如 \olref{fig:counter-Box} 所示，你可以验证它是 $\Box(p \lor q) \lif (\Box p \lor \Box q)$ 的一个反模型。
  \begin{figure}
  \begin{center}
    \begin{tikzpicture}[modal]
      \node[world] (w1) [label={[align=right]right:\mFalse{p}\\ \mFalse{q}}]{$1$};
      \node[world] (w2) [label={[align=right]right:\mFalse{p}\\ \mTrue{q}},
        above left=of w1]{$1.1$};
      \node[world] (w3) [label={[align=right]right:\mTrue{p}\\ \mFalse{q}},
        above right=of w1] {$1.2$};
      \draw[->] (w1) to (w2);
      \draw[->] (w1) to (w3);
    \end{tikzpicture}
  \end{center}
  \caption{$\Box(p \lor q) \lif (\Box p \lor \Box
    q)$ 的一个反模型。}
  \ollabel{fig:counter-Box}
\end{figure}
\end{ex}
}
{
\begin{ex}
  我们知道 $\Proves/ (\Diamond p \land \Diamond q) \lif \Diamond(p
  \land q)$。分析表的构造始于：
  \begin{oltableau}
    [\pFmla{\False}{(\Diamond p \land \Diamond q) \lif \Diamond(p \land q)}{1},
      just = \TAss, checked
      [\pFmla{\True}{\Diamond p \land \Diamond q}{1},
        just = {\TRule{\False}{\lif}[1]}, checked
        [\pFmla{\False}{\Diamond(p \land q)}{1},
          just = {\TRule{\False}{\lif}[1]}
          [\pFmla{\True}{\Diamond p}{1},
            just = {\TRule{\True}{\land}[2]}, checked
            [\pFmla{\True}{\Diamond q}{1},
              just = {\TRule{\True}{\land}[2]}, checked
              [\pFmla{\True}{p}{1.1},
                just = {\TRule{\True}{\Diamond}[4]}, checked
                [\pFmla{\True}{q}{1.2},
                  just = {\TRule{\True}{\Diamond}[5]}, checked
                ]
              ]
            ]
          ]
        ]
      ]
    ]
  \end{oltableau}
  这个分析表尚未完成。下一步，我们考虑唯一没有对勾标记的行：第~$3$ 行上的带前缀公式
  $\sFmla{\True}{\Diamond(p \land q)}[1]$。根据封闭分析表的构造要求，需要对分支上使用的每个前缀应用 $\TRule{\True}{\Diamond}$ 规则，即对 $1.1$ 和 $1.2$ 都要应用：
  \begin{oltableau}
    [\pFmla{\False}{\Diamond(p \land q) \lif (\Diamond p \land \Diamond q)}{1},
      just = \TAss, checked
      [\pFmla{\True}{\Diamond p \land \Diamond q}{1},
        just = {\TRule{\False}{\lif}[1]}, checked
        [\pFmla{\False}{\Diamond (p \land q)}{1},
          just = {\TRule{\False}{\lif}[1]}
          [\pFmla{\True}{\Diamond p}{1},
            just = {\TRule{\True}{\land}[2]}, checked
            [\pFmla{\True}{\Diamond q}{1},
              just = {\TRule{\True}{\land}[2]}, checked
              [\pFmla{\True}{p}{1.1},
                just = {\TRule{\True}{\Diamond}[4]}, checked
                [\pFmla{\True}{q}{1.2},
                  just = {\TRule{\True}{\Diamond}[5]}, checked
                  [\pFmla{\False}{p \land q}{1.1},
                    just = {\TRule{\False}{\Diamond}[3]}
                    [\pFmla{\False}{p \land q}{1.2},
                      just = {\TRule{\False}{\Diamond}[3]}
                    ]
                  ]
                ]
              ]
            ]
          ]
        ]
      ]
    ]
  \end{oltableau}
  现在第 3、8、9 行没有对勾标记。但由于没有添加新的前缀，所以我们对第 8 行和第 9 行应用 $\TRule{\False}{\land}$ 规则，应用于所有由此产生且尚未闭合的分支：
  \begin{oltableau}
    [\pFmla{\False}{(\Diamond p \land \Diamond q) \lif \Diamond(p \land q)}{1},
      just = \TAss, checked
      [\pFmla{\True}{\Diamond p \land \Diamond q}{1},
        just = {\TRule{\False}{\lif}[1]}, checked
        [\pFmla{\False}{\Diamond(p \land q)}{1},
          just = {\TRule{\False}{\lif}[1]}
          [\pFmla{\True}{\Diamond p}{1},
            just = {\TRule{\True}{\land}[2]}, checked
            [\pFmla{\True}{\Diamond q}{1},
              just = {\TRule{\True}{\land}[2]}, checked
              [\pFmla{\True}{p}{1.1},
                just = {\TRule{\True}{\Diamond}[4]}, checked
                [\pFmla{\True}{q}{1.2},
                  just = {\TRule{\True}{\Diamond}[5]}, checked
                  [\pFmla{\False}{p \land q}{1.1},
                    just = {\TRule{\False}{\Diamond}[3]}, checked
                    [\pFmla{\False}{p \land q}{1.2},
                      just = {\TRule{\False}{\Diamond}[3]}, checked
                      [\pFmla{\False}{p}{1.1},
                        just = {\TRule{\False}{\land}[8]}, checked, close
                      ]
                      [\pFmla{\False}{q}{1.1},
                        just = {\TRule{\False}{\land}[8]}, checked
                        [\pFmla{\False}{p}{1.2},
                          just = {\TRule{\False}{\land}[9]}, checked]
                        [\pFmla{\False}{q}{1.2},
                          just = {\TRule{\False}{\land}[9]}, checked, close]
                      ]
                    ]
                  ]
                ]
              ]
            ]
          ]
        ]
      ]
    ]
  \end{oltableau}
  剩下一个开放分支，并且它是完全的。我们据此定义模型：可能世界集为 $W = \{1, 1.1, 1.2\}$
  （即该开放分支上仅出现的前缀），可达关系为 $R = \{\tuple{1, 1.1}, \tuple{1, 1.2}\}$，赋值为 $V(p) =
  \{1.1\}$（因为第~6 行包含 $\sFmla{\True}{p}[1.1]$）和 $V(q) = \{1.2\}$（因为第~7 行包含 $\sFmla{\True}{q}[1.1]$）。该模型如 \olref{fig:counter-Diamond} 所示，你可以验证它是 $(\Diamond p \land \Diamond q) \lif \Diamond (p \land q)$ 的一个反模型。
  \begin{figure}
  \begin{center}
    \begin{tikzpicture}[modal]
      \node[world] (w1) [label={[align=right]right:\mFalse{p}\\ \mFalse{q}}]{$1$};
      \node[world] (w2) [label={[align=right]right:\mTrue{p}\\ \mFalse{q}},
        above left=of w1]{$1.1$};
      \node[world] (w3) [label={[align=right]right:\mFalse{p}\\ \mTrue{q}},
        above right=of w1] {$1.2$};
      \draw[->] (w1) to (w2);
      \draw[->] (w1) to (w3);
    \end{tikzpicture}
  \end{center}
  \caption{$(\Diamond p \land \Diamond q) \lif
    \Diamond (p \land q)$ 的一个反模型。}
  \ollabel{fig:counter-Diamond}
\end{figure}
\end{ex}
}

\end{document}
